\scanpage{304}
\mspage{3}
\[ {K^{[m]}}_{[n]}=K(\Delta^m,\Delta^n)\in A \]
De \unsure{m\textsuperscript{e}}, on note
\[
\begin{aligned}
{K_*}^*&\quad\text{les objets de }(A^*)_*\simeq\underline{\operatorname{Hom}}(\boldsymbol\Delta^o\times\boldsymbol\Delta,A)\\
K^{**}&\quad\text{les objets de }(A^*)^*\simeq\underline{\operatorname{Hom}}(\unsure{\boldsymbol\Delta^o|\boldsymbol\Delta}\times\boldsymbol\Delta,A)\\
K_{**}&\quad\text{les objets de }A_{**}=\underline{\operatorname{Hom}}(\boldsymbol\Delta^o\times\boldsymbol\Delta^o,A)
\end{aligned}
\]

On fera attention que si on compose avec l'un des foncteurs de symétrie
\deleted{$(m,n)\longmapsto(n,m)$}
\inserted{$(\Delta^m,\Delta^n)\longmapsto(\Delta^n,\Delta^m)$}
dans $\boldsymbol\Delta\times\boldsymbol\Delta$ ou dans $\boldsymbol\Delta^o\times\boldsymbol\Delta^o$, ou entre $\boldsymbol\Delta\times\boldsymbol\Delta^o$ et $\boldsymbol\Delta^o\times\boldsymbol\Delta$), on n'utilisera pas la même lettre pour désigner un \deleted{complexe} \inserted{objet} (tel que ${K^*}_*$ ou $K^{**}$ disons) ou son transformé (qui devrait alors s'écrire ${K_*}^*$, ou encore une fois $K^{**}$), car cela conflicterait avec la convention plus haut de la ``montée et descente des indices''
\[ {K_*}^*\simeq({K^*}_*)^\vee \]
dans le premier cas, et \deleted{\unsure{serait}} amènerait à écrire
\[ K^{[m][n]}=K^{[n][m]} \]
dans le deuxième ! On utilisera au besoin un symbole ``de symétrie'' tel que $\sigma$, pour dénoter le symétrique de $K$ ($={K^*}_*,K^{**}$) par
\[ {}^\sigma K\ (= {({}^\sigma K)_*}^*,({}^\sigma K)^{**}) \]
\deleted{\unsure{de}} \ldots D'ailleurs,
\deleted{Ces considérations et notations}
si on a un objet ${K^*}_*$ disons, \inserted{$\in\operatorname{Ob}(\unsure{A_*})^*$} alors pour $\forall n\geq0$, on peut considérer
\[
\deleted{\unsure{u\longmapsto}}\quad
{K^*}_{[n]}\overset{\mathrm{def}}{=}(m\longmapsto {K^{[m]}}_{[n]})
\]
comme un objet de $A^*$ ; le foncteur $[n]\longmapsto {K^*}_{[n]}$ de ${A^*}_*$ n'est autre que ${}^\sigma K$.
On étend ces considérations et notations au cas de symboles $K$ à plus de deux indices, pour désigner des objets de catégories de

\margin{305}
